% Please do not change %
\documentclass[11pt]{article} % font size
\usepackage[margin=1in]{geometry} % margins
\usepackage{amsmath,amsthm,amssymb} % for the  common "math symbols"
\usepackage[shortlabels]{enumitem} % for enumeration
\usepackage{booktabs} % if you need tables
\usepackage{listings}
\usepackage{color}

\newtheorem{prop}{Proposition}

\lstset{
  frame=none,
  xleftmargin=2pt,
  stepnumber=1,
  numbers=left,
  numbersep=5pt,
  numberstyle=\ttfamily\tiny\color[gray]{0.3},
  belowcaptionskip=\bigskipamount,
  captionpos=b,
  escapeinside={*'}{'*},
  language=haskell,
  tabsize=2,
  emphstyle={\bf},
  commentstyle=\it,
  stringstyle=\mdseries\rmfamily,
  showspaces=false,
  keywordstyle=\bfseries\rmfamily,
  columns=flexible,
  basicstyle=\small\sffamily,
  showstringspaces=false,
  morecomment=[l]\%,
}
% Feel free to include additional packages (if needed), but make sure it does not override the margin/font requirements.
%
\setlength{\parindent}{0pt} % no indent for new paragraphs.
\setlength{\parskip}{5pt} % distance between paragraphs, which will be used when you have a blank line in your LaTeX code.


%%%%%%%%%%%%%%%%%%%%%%%%%%%%%%%%%%%%%%%%%%%%%%%%%%%%%%%%%%%%%%%%%%%%%%%%%%%%%%
\title{Principles of Programming Languages: \\Problem Sheet 4 - Continuations and types}
%
\author{Wira Azmoon Ahmad}
\date{27 Nov 2023}
%
\begin{document}
%
\maketitle

%You may use the usual \LaTeX\ environments to structure your answers:
%\begin{align}
%        \label{basegnn}
%        \mathbf{h}_u^{(t+1)} = \sigma \bigg( \mathbf{W}_\text{self}^{(t)} \mathbf{h}_u^{(t)}  + \mathbf{W}_\text{neigh}^{(t)} \sum_{v \in N(u)} \mathbf{h}_v^{(t)} + \mathbf{b}^{(t)} \bigg). 
%\end{align}


%\begin{table}[h]
%\caption{A simple table.}
%\centering
%\begin{tabular}[t]{lrrrr}
%\toprule
%Graphs & Can & Be  & Fun & Too \\ 
%\midrule 
%$G_1$ & $1$  & $2$ & $3$ & $4$ \\ 
%$G_2$ & $-1$  & $-2$ & $-3$ & $-4$ \\ 
%$G_3$ & $0.1$  & $0.2$ & $0.3$ & $0.4$ \\ 
%\bottomrule
%\end{tabular}
%\label{tab}
%\end{table}

\section*{Question 1}

% To answer each part of the question, please simply use a paragraph command 
% \paragraph{(a)}


\lstinputlisting[language=Haskell]{../revk.hs}
Suppose we have a list $[a_1,a_2,...,a_n]$.
We show by induction that the continuation for iteration $i$ can be written in the form
\[(...((y_i ++ [a_{i-1}]) ++ [a_{i-2}])++...++[a_1]) = y_i + [a_{i-1},...,a_1] \]

\lstinputlisting[language=Haskell, lastline=3]{../PS4.hs}

Suppose for $n=k$, the continuation is 
\[y_k + [a_{k-1},...,a_1] \]

\lstinputlisting[language=Haskell, linerange={5-8}]{../PS4.hs}
We write in \texttt{fun}:
\lstinputlisting[language=Haskell]{../revk.fun}

\section*{Question 2}
\lstinputlisting[language=Haskell]{../fibcpsdefun.hs}
\lstinputlisting[language=Haskell,linerange={10-41}]{../PS4.hs}

\section*{Question 3}
Charlie is right because the point of CPS is to prevent unbounded recursion, 
and the division here has no recursion.

\section*{Question 4}
\paragraph{(a)}
\begin{prop}
$\forall s,s\in S \implies \exists n\in\mathbb{Z}_{\geq 0} : s = ab^n$
\end{prop}
\begin{proof}
  Let $P(s)=(\exists n\in\mathbb{Z}_{\geq 0} : s = ab^n)$. The first rule is true, since for $n=1$, $ab\in S$.

  For rule 2, suppose $P(ax)$. Then $\exists n\in\mathbb{Z}_{\geq 0} : ax = ab^n$. So $axx = ab^{2n}$.
  Let $n'=2n\in\mathbb{Z}_{\geq 0}$, then $\exists n'\in\mathbb{Z}_{\geq 0} : axx = ab^{n'}$. Then $P(axx)$.

  For rule 3, suppose $P(abbbx)$. Then $\exists n\in\mathbb{Z}_{\geq 0} : abbbx = ab^n$, where $n\geq3$. 
  So $ax=ab^{n-3}$. 
  Let $n'=n-3\in\mathbb{Z}_{\geq 0}$, then $\exists n'\in\mathbb{Z}_{\geq 0} : ax = ab^{n'}$. Then $P(ax)$.
\end{proof}

\paragraph{(b)}
\[\dfrac{
  \dfrac{
    \dfrac{
      \dfrac{
        \dfrac{-}{ab\in S}(1)
      }{abb\in S} (2), x=b
    }{abbbb \in S} (2), x=bb
  }{abbbbbbbb \in S} (2), x = bbbb
}{abbbbbbbb \in S}(3), x = bbbbb\]


\begin{prop}
$\forall s,s\in S \implies \exists n\in \{k:k\in \mathbb{Z}^+ \wedge 3 \nmid k\} : s = ab^n$.
In other words, $n$ is not a multiple of 3.
\end{prop}

\begin{proof}
  Let $T=\{k:k\in \mathbb{Z}^+ \wedge 3 \nmid k\}$.

  Let $P(s)=(\exists n\in T : s = ab^n)$. The first rule is true, since for $n=1\in T$, $ab\in S$.

  For rule 2, suppose $P(ax)$. Then $\exists n\in T : ax = ab^n$. So $axx = ab^{2n}$.
  Let $n'=2n\in T$, then $\exists n'\in T : axx = ab^{n'}$. Then $P(axx)$.

  For rule 3, suppose $P(abbbx)$. Then $\exists n\in T : abbbx = ab^n$, where $n>3$. 
  So $ax=ab^{n-3}$. Since $3\nmid n$, $3 \nmid (n-3)$.
  Let $n'=n-3\in T$, then $\exists n'\in T : ax = ab^{n'}$. Then $P(ax)$.
\end{proof}

\clearpage
\section*{Question 5}
On top of the transition rules in lecture, we have the following
\begin{equation}
\begin{aligned}
&\langle succ, env, k \rangle \rightarrow [k, succ]\\
&\langle pred, env, k \rangle \rightarrow [k, pred]\\
&[succ(\Box):k,v] \rightarrow \langle v+1, -, k\rangle\\
&[pred(\Box):k,v] \rightarrow \langle v-1, -, k\rangle\\
\end{aligned}
\end{equation}
\begin{equation}
\begin{aligned}
  &\langle pred (succ (1)), -, Show \rangle \\
  &\rightarrow \langle pred, -, (\Box(succ(1)), -):Show\rangle\\
  &\rightarrow [ (\Box(succ(1)), -):Show, pred]\\
  &\rightarrow \langle succ(1), -, pred(\Box): Show \rangle\\
  &\rightarrow \langle succ, -, (\Box(1), -):pred(\Box): Show \rangle\\
  &\rightarrow [ (\Box(1), -):pred(\Box): Show, succ ]\\
  &\rightarrow \langle 1, -, succ(\Box) :pred(\Box): Show \rangle\\
  &\rightarrow [ succ(\Box) :pred(\Box): Show , 1]\\
  &\rightarrow \langle 2, -, pred(\Box): Show \rangle\\
  &\rightarrow [ pred(\Box) : Show , 2]\\
  &\rightarrow \langle 1, -, Show \rangle
\end{aligned}
\end{equation}

\clearpage
\section*{Question 6}
\paragraph{(a)}
No to both.
Consider the infinite chain 
\[1\leq2\leq3\leq...\]
which clearly does not have any upper bound.
It has no bottom element since for any integer $x$, we can also find an integer $y=x-1$ so $x\not\leq y$.

\paragraph{(b)}
Take the empty set $\emptyset\subset\mathbb{Z}$.
It is vacuously true that $\forall y \in \mathbb{Z}, \forall x\in \emptyset,y\leq x$,
so all $y\in\mathbb{Z}$ are upper bounds for any chain from $\emptyset$,
and since $\mathbb{Z}$ has no least element, there is no least upper bound.

\paragraph{(c)}
\begin{prop}
  For any chain $S_1\subseteq S_2 \subseteq ... \subseteq S_n \subseteq ...$,
  the set $S = \bigcup_i S_i$ is the least upper bound.
\end{prop}
\begin{proof}
  Clearly, $\forall i, S_i \subseteq S$. So $S$ is an upper bound.

  Suppose we have another upper bound $S_y$. 
  As it is an upper bound, $\forall i, S_i \subseteq S_y$.
  Then $\bigcup_i S_i \subseteq S_y$.
  By definition, $S = \bigcup_i S_i \subseteq S_y$. So $S$ is the least upper bound.
\end{proof}

The bottom element is $\perp = \emptyset$ since $\forall S_i\in P(\mathbb{N}), \emptyset \subseteq S_i$.

\begin{prop}
  $F(S)=\{1\} \cup \{n+2|n\in S\}$ is continuous.
\end{prop}
\begin{proof}
  Suppose $S_x \subseteq S_{x'}$.
  Then $F(S_x) = \{1\} \cup \{n+2|n\in S_x\} \subseteq \{1\} \cup \{n+2|n\in S_{x'}\} = F(S_{x'})$.
  So $F$ is monotone.

  Suppose we have a chain $S_1\subseteq S_2 \subseteq ...$.
  We show that $F(\bigcup_i S_i) = \bigcup_i F(S_i)$.
  \begin{equation}
  \begin{aligned}
    F(\bigcup_i S_i) &= \{1\} \cup \{n+2|n\in \bigcup_i S_i\}\\
    &= \{1\} \cup \bigcup_i \{n+2|n\in  S_i\}\\
    &= \bigcup_i \{1\} \cup  \{n+2|n\in  S_i\}\\
    &= \bigcup_i F(S_i)
  \end{aligned}
  \end{equation}
  So $F$ preserves least upper bounds of chains.

  Then $F$ is continuous.

\end{proof}

By Tarski's fixed point theorem, the least fixed point of $F$ is
\[\bigcup_n F^n(\emptyset)=\{k|2k-1\in\mathbb{N}\}\]
In other words, it is set of odd natural numbers.

\end{document}